\subsection{Feasibility}
\vspace{12pt}

The proposed research on blockchain scalability is highly feasible, drawing upon a robust foundation of existing technologies, ongoing industry developments, and a strong market demand for enhanced blockchain performance.

\subsubsection{Technical Feasibility}
The extensive literature reviewed demonstrates that various technical approaches to blockchain scalability---including on-chain modifications (e.g., sharding, DAGs, block compression), off-chain solutions (e.g., payment channels, sidechains, cross-chains), advanced consensus mechanisms, and network/platform-layer optimisations---have been proposed, prototyped, and in some cases partially implemented. Frameworks such as Ethereum~2.0 and architectural advancements in Hyperledger Fabric exemplify the active pursuit of scalable blockchain solutions. These diverse methods confirm that improving blockchain performance is technically achievable, albeit complex, involving trade-offs. The feasibility of this research relies on synthesising these existing technical efforts rather than proposing entirely new, unproven technologies.

\subsubsection{Resource Availability}
Conducting a systematic review requires access to academic databases and literature management tools. The specified databases (IEEE Xplore, ACM Digital Library, ScienceDirect, Springer Link, Web of Science, Google Scholar, Scopus) are widely accessible within academic institutions. Tools like EndNote or Zotero are standard for managing references and screening papers, ensuring efficient execution of the methodology. Thus, the necessary informational and software resources are readily available.

\subsubsection{Data Availability}
The empirical performance metrics and analyses, as presented in the reviewed research, provide a rich source of data for comparative evaluation. Studies on throughput, latency, and storage, as well as qualitative assessments of security and decentralisation, are well documented. This ensures feasibility of the proposed analysis without requiring new data generation from scratch. These align closely with technical KPIs such as \textit{Transactions per Second (TPS)}, \textit{Transaction Latency}, \textit{Block Size/Time}, and \textit{Network Uptime}:contentReference[oaicite:2]{index=2}.

\subsubsection{Expertise}
The research builds upon established concepts in distributed ledger technologies, cryptography, computer networks, and performance analysis. Researchers with backgrounds in these fields---such as the authors of the provided sources (e.g., Sanka and Cheung, with expertise in blockchain technology, digital systems design, network computing, and information security)---possess the requisite knowledge to undertake such a comprehensive study.

\subsubsection{Societal and Economic Feasibility}
There is a significant and growing demand for scalable blockchain solutions across industries. The forecasted business value of blockchain (over \$3.1 trillion by 2030) and the high percentage of companies planning to integrate it underscore the economic incentive. Addressing scalability is crucial for blockchain to transition from experimental stages to mainstream production, enabling adoption in decentralised finance (DeFi), supply chain, healthcare, and IoT. This validates both the societal and economic relevance of the research.

\subsubsection{Current Adoption Trends}
The increasing adoption of blockchain by major companies (e.g., Microsoft, IBM, Oracle, Ford, BMW, Renault, Maersk) and countries (e.g., El-Salvador, Georgia, UAE, Singapore) for diverse applications demonstrates real-world demand. This ongoing integration across sectors indicates a sufficiently mature ecosystem to support further scalability research, making the study timely and impactful.

\subsection{Conclusion}
In conclusion, the proposed research is highly feasible, supported by existing technological advancements, readily available academic resources, a wealth of empirical data, and a clear market demand. By systematically synthesising state-of-the-art solutions and evaluating them against well-defined KPIs---spanning \textit{technical performance}, \textit{operational efficiency}, \textit{economic viability}, and \textit{security}:contentReference[oaicite:3]{index=3}:contentReference[oaicite:4]{index=4}---this study aims to provide a comprehensive understanding of blockchain scalability. The findings will contribute both theoretically and practically, guiding the design and adoption of scalable blockchain systems.



\vspace{12pt}
\subsection{Conclusion}
\vspace{12pt}

This project proposal outlines an innovative approach to optimize container startup times in Kubernetes environments through self-adaptive CPU allocation. By leveraging predictive algorithms and integrating with the Kube Startup CPU Boost feature, the project aims to significantly reduce application startup times and enhance resource allocation efficiency. The proposed solution addresses the critical challenge of efficient resource allocation during the resource-intensive startup phase of applications, combining real-time monitoring, predictive analysis, and startup boost to ensure optimal CPU allocation without over/under provisioning.

\vspace{12pt}

The potential impact of this project is substantial, with expected improvements in application performance, reduced operational/management costs, and enhanced overall efficiency of Kubernetes clusters. The comprehensive evaluation plan, including a range of benchmarks and performance metrics, will provide a thorough assessment of the solution's effectiveness. Given its technical viability and potential economic benefits, this project represents a significant step forward in container orchestration and resource management, with the potential to contribute valuable insights and practical solutions to the cloud computing community.

\vspace{12pt}